% !PRINTED 441
\weberpage{441}
% running head: § 121. Trennung der entgegengesetzten Klassen.

Ist $k$ die \gesperrt{Hauptklasse}, so ist
\[
2\omega = \sqrt{D} \text{ oder } 2\omega = -1+\sqrt{D},
\]
und da $\sqrt{D}$ rein imaginär ist, so ist $k=j(\omega)$ nach § 69 (4) \gesperrt{reell}. % [sic] Der Druck schreibt k=j(\omega), obwohl die eingeführte Klasseninvariante mit (k)=j(\omega) bezeichnet wird.
Demnach ergibt sich aus (14) und (15), da jetzt $(sk)=(s)$ wird:

\gesperrt{Die Invarianten entgegengesetzter Klassen sind konjugiert imaginär, die Invarianten zweiseitiger Klassen sind reell.}

Kehren wir zu einer beliebigen Klasse $k$ zurück und setzen in der Formel (15) $sk$ an Stelle von $k$, so folgt:
\weq{16}{(k) = f_s(sk, -\sqrt{D}),}
oder, indem man $sk=k'$ setzt und $f$ für $f_s$ schreibt, nach (21) % [sic] printed (21); this is a forward reference in a passage currently numbered (16)--(17)
\weq{17}{\begin{aligned} (k') &= f(k, \sqrt{D}),\\ (k) &= f(k', -\sqrt{D}), \end{aligned}}
worin nun $k, k'$ irgend zwei Klassen der Determinante $D$ sein können.

Hiernach sind wir imstande, die \gesperrt{Galoissche} Gruppe der Klassengleichung, oder zunächst wenigstens eine Gruppe von Permutationen zu bestimmen, in der die Gruppe der Klassengleichung als Teiler enthalten ist.

Nehmen wir zunächst an, es sei die $\sqrt{D}$ dem Rationalitätsbereich der rationalen Zahlen adjungiert, und wenn $(k), (k'), (k'')\ldots$ die sämtlichen Klasseninvarianten der Diskriminante $D$ bedeuten,
\[
R(k, k', k'', \ldots)
\]
irgend eine rationale Funktion dieser Größen. Nach unserem Satze läßt sich diese Funktion rational ausdrücken durch eine der Größen $(k)$, also etwa:
\[
R(k, k', k'', \ldots) = R'(k).
\]
Bedeutet ferner $(s)$ eine beliebige Klasseninvariante der Diskriminante $D$, und hat die Funktion $R$ die Eigenschaft, durch die sämtlichen $h$ Permutationen
\[
\begin{pmatrix} k, & k', & k'', & \ldots \\ sk, & sk', & sk'', & \ldots \end{pmatrix},
\]
deren Gesamtheit wir mit $\mathfrak{S}$ bezeichnen wollen, ungeändert zu bleiben, so ist $R'(k)=R'(sk)$, und daher \gesperrt{rational ausdrückbar}.

\gesperrt{Die Galoissche Gruppe der Klassengleichung nach Adjunktion von $\sqrt{D}$ ist also in dem System $\mathfrak{S}$ enthalten,}
%CONT>>
